% !TeX program = xelatex
% UTF-8; standalone XeLaTeX/XePersian document.
% When included in an existing XePersian document, the preamble is skipped.
% The earlier five-page transcription is retained verbatim below.
% Original directory dermatology/atopic dermatitis is preserved unchanged.
% New source: Google Drive Notebook/atopic dermatitis, filename order:
% images/pic01.jpg = 20261018_200002688.jpg; 1qxFUb2a9yfgRUp5TpASCyheqcjqagQ_d
% images/pic02.jpg = 20261018_200006792.jpg; 1DQw2TfhqTR20AmYnHFf2ry2_iTc8LIll
% images/pic03.jpg = 20261018_200013991.jpg; 1j8a-TVzDuOWAHZt6lMdRsRmLGKrjcsis
% images/pic04.jpg = 20261018_200018908.jpg; 1rEMC7gxesBq7dUCPy1fQJHkm3fHceYvr
% Compile from this directory: xelatex "atopic dermatitis.tex"
% Packages graphicx, amsmath, amssymb, xcolor, enumitem must precede
% xepersian in a host document that inputs this file.
\ifdefined\settextfont
  \let\PADfinish\relax
\else
  \def\PADfinish{\end{document}}
  \documentclass[12pt,a4paper]{article}
  \usepackage[margin=2cm]{geometry}
  \usepackage{amsmath,amssymb,xcolor,enumitem,graphicx}
  \usepackage{xepersian}
  \IfFontExistsTF{Amiri}{\settextfont{Amiri}}{\settextfont{FreeSerif}}
  \setlatintextfont{FreeSerif}
  \setlength{\parindent}{0pt}
  \setlength{\parskip}{5pt}
  \setlength{\emergencystretch}{3em}
  \begin{document}
\fi
\begingroup
\begin{center}
{\Large\bfseries درماتیت آتوپیک}\par
{\large یادداشت‌های پیشین و چهار تصویر جدید}
\end{center}
\section*{بخش نخست: یادداشت‌های پیشین ـ اگزما آتوپیک ـ ۴۱}
% Earlier transcription: all text, including its uncertainty markers, retained.
\providecommand{\en}[1]{\lr{#1}}
\providecommand{\adunclear}[1]{\textcolor{red!65!black}{[خوانش نامطمئن: #1]}}
\providecommand{\adcut}{\textcolor{red!65!black}{[ادامه خارج از کادر عکس]}}
\providecommand{\adpage}[1]{\par\bigskip\noindent\textbf{صفحهٔ #1 ـ اگزما آتوپیک ـ ۴۱}\par\smallskip\noindent موضوع: اگزما آتوپیک ـ ۴۱ \hfill تاریخ: / /\par\medskip}

\adpage{۱}
\noindent\en{Rook's Dermatology (10th)}

یک وضعیت خارش‌دار، عودکننده، به‌طور مزمن و التهابی پوست است که عموماً در اوایل کودکی آغاز می‌شود (معمولاً قبل از ۲ سالگی).

\begin{itemize}
  \item ویژگی بارز: التهاب (قرمزی در پوست‌های روشن و افزایش التهاب در پوست‌های \adcut{} تیره).
  \item تغییرات سطح پوست: \en{Fine scaling}، اکسکوریشن، لیکنیفیکاسیون.
  \item معمولاً در زمینهٔ خارش مداوم؛ ادم، پاپولیشن، اوزینگ (\adunclear{تراوش؛ انتهای سطر در لبهٔ عکس}).
  \item عموماً سطوح فلکسورال را درگیر می‌کند.
  \begin{itemize}
    \item درگیری غیر فلکسورال در شیرخواران شایع است و سیاه‌پوستان \adcut
  \end{itemize}
\end{itemize}

\en{AD} معمولاً با سایر وضعیت‌های آتوپیک همراهی دارد.

\subsection*{تظاهرات بالینی}
خارش، اریتم ماکولار، پاپول/پاپولووزیکول، ضایعات اگزماکوز ـ به‌ویژه \adunclear{کراتوتیک}، لیکنیفیکاسیون، هایپر و هایپوپیگمانتاسیون \adcut، خشکی پوست، عفونت ثانویه.

\subsection*{کرایتریا تشخیص}
کودک باید حتماً پوست خارش‌دار طی ۱۲ ماه گذشته \adcut؛ حداقل ۳ مورد:
\begin{enumerate}
  \item سن شروع زیر ۲ سال.
  \item سابقهٔ درگیری فلکسورال.
  \item سابقهٔ پوست خشک.
  \item سابقهٔ سایر بیماری‌های آتوپیک.
  \item درماتیت فلکسورال قابل مشاهده.
\end{enumerate}
\noindent حاشیهٔ معیار پنجم: آسم در سن زیر ۴ سال \adcut؛ اقوام درجهٔ ۱ هم حساب \adcut

\adpage{۲}
\begin{itemize}
  \item بیماران \en{AD} می‌توانند سطح نرمال \en{total IgE} داشته باشند.
  \item \en{AD} \adunclear{لزوماً} مرتبط با یک آلرژن خاص است.
\end{itemize}

\begin{latin}
\noindent AD $\longrightarrow$ Allergen-specific IgE $\uparrow$ (40--90\% of cases)\\
\hspace*{2em}(Extrinsic AD)\par\medskip
\noindent AD $\longrightarrow$ N/L IgE level (50\% $<$ of cases)\\
\hspace*{2em}(Intrinsic AD or non-atopic AD)
\end{latin}
% The comparison sign above follows the handwritten English line literally.

\begin{itemize}
  \item در طول زمان و افزایش عمر موارد \en{E-AD} شایع‌تر می‌شوند (تا ۹۵٪).
  \item شیوع \en{AD} در ایران ۲٪ و در ژاپن ۱۶٪ است؛ \en{$M<F$}.
  \item شیوع سالانهٔ آن حدود ۲۰٪ است.
  \item شیوع آن در میان‌سالی و پیری نیز افزایش می‌یابد.
  \item در \en{AD}: عوامل ژنتیکی و عوامل محیطی $\leftarrow$ تغییر عملکرد پوست $\leftarrow$ التهاب پوست $\leftarrow$ \en{AD}.
  \item موارد شدید \en{AD} نامعمول هستند و تنها ۵٪ موارد را تشکیل می‌دهند.
  \item ۲–۱ وعده میوهٔ تازه در هفته احتمال \en{AD} را کم می‌کند.
  \item ۳ بار فست‌فود در هفته احتمال \en{AD} را زیاد می‌کند.
\end{itemize}

\adpage{۳}
\begin{itemize}
  \item مصرف ماهی طی بارداری ریسک \en{AD} طی ۵ سال اول زندگی را \adunclear{کم می‌کند؛ پایان سطر بریده است}.
  \item مکمل‌های تغذیه‌ای نقش خاصی ندارند.
  \item تغذیه با شیر مادر احتمال \en{AD} را کاهش می‌دهد.
  \item چاقی هم ریسک بروز و هم شدت \en{AD} را افزایش می‌دهد.
  \item مصرف دخانیات در پدر و مادر ریسک \en{AD} را تغییر نمی‌دهد.
  \item آلودگی هوا به \en{NO\textsubscript{2}} تغییر نمی‌دهد.
  \item مهدکودک در ۲ سال اول زندگی ریسک را بالا می‌برد.
  \item حیوان خانگی در سال‌های اولیهٔ زندگی ریسک را کم می‌کند (خصوصاً سگ).
\end{itemize}

\subsection*{مارش آتوپیک}
\begin{itemize}
  \item پیشروی \en{AD} به سایر بیماری‌های آتوپیک:
  \begin{itemize}
    \item رینیت آلرژیک، آلرژی غذایی، آسم.
    \item تنها در ۷٪ موارد \en{AD} رخ می‌دهد.
  \end{itemize}
  \item \en{IgE} نقش جزئی در \en{AD} داشته چرا که مهار آن کارساز نیست $\leftarrow$ اومالیزومب.
  \item در \en{AD} سطح ائوزینوفیل و نوتروفیل بالاست.
  \item ۲۱٪ بیماران \en{AD} آلرژی غذایی دارند (۴٪ در نرمال جمعیت).
  \item ارتباطی بین درماتیت آلرژیک تماسی (\en{ACD}) و \en{AD} وجود ندارد.
\end{itemize}

\adpage{۴}
\begin{itemize}
  \item خارش در \en{AD}: سیگنالینگ غیر هیستامینرژیک دارد.
  \begin{itemize}
    \item با جزء التهابی \en{AD} مرتبط است: \en{IL-31}، \en{TSLP}.
    \item به همین جهت با داروهای ضدالتهابی رفع می‌شود.
    \item سیکلوسپورین و مهارکننده‌های \en{JAK}.
    \item سطح \en{IL-31} با شدت خارش مرتبط است.
    \item سطح \en{IL-31} به دنبال سوپرآنتی‌ژن استافیلوکوک به‌شدت بالا می‌رود.
  \end{itemize}
\end{itemize}

\subsection*{تظاهرات بالینی \en{AD}}
\begin{itemize}
  \item سن شایع زیر ۲ سال است؛ خصوصاً ۶–۲ ماهگی.
  \item علامت اصلی خارش، به همراه راش مزمن \en{fluctuating} است.
  \item اختلال خواب به علت خارش؛ اضطراب و تخریب پوستی.
  \item محرک خارش: گرما، عرق، حمام، ورزش، استرس، لباس پشمی.
  \item ویژگی ضایعهٔ اگزما حاد: \adunclear{سنتریپتال}، حاشیه نامنظم، ملتهب، اندکی برجسته؛ بیشتر سطوح فلکسورال.
  \item اگزما مزمن: پوست کلفت می‌شود، لیکنیفیکاسیون، هایپرپیگمانتاسیون.
  \item کراست‌های زرد $\leftarrow$ عفونت ثانویهٔ باکتریال.
  \item پوست خشک است به دلیل تبخیر زیاد، به دلیل اختلال ساختاری پوست.
  \item در شیرخواران: بیشتر در صورت / سطوح فلکسورال.
\end{itemize}

\adpage{۵}
\begin{itemize}
  \item در کودکان: بیشتر سطح فلکسور، آرنج و زانو / سطوح کناری گردن / مچ دست و مچ پا.
  \item پچ‌های \adunclear{دیسکرومیه} ممکن است دیده شود.
  \item درگیری دست که گاهی با وزیکولیشن همراه است ممکن است دیده شود (\en{Pompholyx}).
  \item اتواینوکولیشن \adunclear{دلیل اول} عفونت باکتری و ویروسی را \adcut
  \item حاشیهٔ سطر قبل: \adunclear{بیشتر ...؛ ادامه در لبهٔ چپ عکس بریده است}.
  \item در بزرگسالان:
  \begin{itemize}
    \item بیشتر سطوح فلکسورها و دست.
    \item لیکنیفیکاسیون شایع‌تر است.
  \end{itemize}
\end{itemize}

\subsection*{منج درماتیت آتوپیک}
\begin{itemize}
  \item باید براساس شدت بیماری و نیاز بیمار باشد.
  \item حتماً شدت بیماری قبل شروع درمان ارزیابی و مکتوب گردد.
  \item درمان خط اول:
  \begin{itemize}
    \item کاهش مواجهه با تحریک‌کننده‌ها.
    \item حمام با مواد مرطوب‌کننده.
    \item استفادهٔ مداوم از مرطوب‌کننده $\leftarrow$ حتی در رمیژن / در همهٔ بیماران \adcut
    \item سرکوب التهاب با کمک \en{TCS} $\leftarrow$ \adunclear{در مرحلهٔ شدید در بیماران هم اتفاق} \adcut
  \end{itemize}
  \item معمولاً نیازی به تغییر رژیم غذایی نیست.
  \item آلرژن‌های تنفسی هم ممکن است دخیل باشند.
\end{itemize}


\clearpage
\section*{بخش دوم: درماتیت آتوپیک ـ ۸۰}
\textbf{یادداشت رونویسی:} متن فارسی و انگلیسی، اعداد، اختصارات و تأکیدها از چهار تصویر جدید بازنویسی شده‌اند. روابطِ فلش‌ها با فهرست و توضیح حفظ شده‌اند. خوانش نامطمئن صریحاً مشخص شده است. مطالب پزشکی منبع بدون اصلاح محتوایی نقل شده‌اند. تصاویر کامل هر دو مجموعه در پیوست آمده‌اند.

\newcommand{\PADrange}[2]{\lr{#1--#2}}
\newcommand{\PADunclear}[1]{\textcolor{red!65!black}{[خوانش نامطمئن: #1]}}
\newcommand{\PADheader}[2]{\section*{تصویر جدید #1}\textbf{سربرگ و حاشیه:}
\lr{Subject:} #2؛ \lr{Year:}، \lr{Month:}، \lr{Date:} (خالی)؛
شماره‌های چاپی کنار سطرها: ۱، ۲، ۳، ۴، ۵، ۶، ۷، ۸، ۹؛ شمارهٔ دست‌نویس پایین صفحه: #1.\par}

\PADheader{۱}{۸۰ ـ درماتیت آتوپیک}
\begin{itemize}
\item \lr{AD}: التهاب، مزمن، \textbf{خارش‌دار}، \textbf{عودکننده}.
\item سلول‌های لانگرهانس فعال‌شده در درم، \lr{IgE} بیان می‌کنند؛ سلول‌های \lr{T} تحریک می‌شوند:
\begin{enumerate}
\item در فرایند حاد، \lr{Th2} فعال‌شده وارد درم می‌شود.
\begin{itemize}
\item \lr{IL-4}، \lr{IL-13}: تولید \lr{IgE}.
\item \lr{IL-5}: تمایز ائوزینوفیل‌ها.
\end{itemize}
\item فرایند مزمن: \lr{Th0} و \lr{Th1}؛ \lr{IL-12}، \lr{IL-18}.
\end{enumerate}
\item تریگرهای بیماری: استرس، دما و رطوبت خیلی کم یا خیلی زیاد، محرک‌ها، آلرژن‌ها و عفونت‌ها.
\item آلرژی به تخم‌مرغ:
\begin{latin}
The most common cause of food-induced eczematous reactions
\end{latin}
\end{itemize}

\subsection*{اپیدمیولوژی}
\begin{itemize}
\item \PADrange{15}{20}\% کودکان مبتلا به \lr{AD} هستند.
\item معمولاً در شیرخوارگی شروع می‌شود.
\item ۵۰\% کودکان مبتلا، ظرف سال اول علامت‌دار می‌شوند؛ \textbf{۸۰\% قبل از ۵ سالگی}.
\item \textbf{\lr{AD} معمولاً اولین تظاهر مارچ آتوپی است}؛ آسم، رینیت آلرژیک؛ «۸۰\% گام بعدی».
\end{itemize}

\subsection*{تظاهرات بالینی}
\begin{itemize}
\item در شیرخواران: \textbf{سطوح اکستانسور اندام‌ها}، \lr{face, scalp, cheeks}.
\item \textbf{ناحیهٔ پوشک درگیر نمی‌شود}.
\item در کودکان با سن بالاتر: راش در ناحیهٔ \textbf{آنتی‌کوبیتال} و پوپلیتئال فوسا، سر و گردن.
\item در نوجوانان و بزرگسالان: \textbf{پلاک‌های لیکنیفیه} در نواحی فلکسور.
\end{itemize}

\clearpage
\PADheader{۲}{۸۰ ـ \lr{AD}}
\begin{itemize}
\item خارش: \textbf{در شب شدیدتر است}؛ اختلال خواب.
\item \lr{Ph/E}: \textbf{درماتوگرافیسم}، هایپرلینیاریتی کف دست‌ها، پیتریازیس آلبا،
\begin{latin}
\textbf{Dennie-Morgan folds} (creases under the lower eyelids).

Keratosis pilaris $\longrightarrow$ Asymptomatic, horny follicular papules,\\
Extensor surfaces of the arms
\end{latin}
\end{itemize}

\subsection*{تشخیص}
\begin{itemize}
\item \textbf{تشخیص \lr{AD} بالینی است}.
\item بیوپسی برای \lr{R/O} کردن علل دیگر ممکن است انجام شود.
\item برای یافتن آلرژن: \lr{Skin Testing, Serum Specific IgE testing}.
\end{itemize}

\subsection*{تشخیص‌های افتراقی و نقص ایمنی}
\begin{itemize}
\item شیرخوار با \lr{FTT}، عفونت راجعهٔ پوستی و سیستمیک، راش اریتماتو: از نظر \textbf{نقایص ایمنی} بررسی شود.
\item سندرم ویسکوت ـ آلدریچ:
\begin{itemize}
\item \textbf{\lr{X-linked recessive}}.
\item \textbf{\lr{AD} + ترومبوسایتوپنی + پلاکت‌های کوچک}.
\item \textbf{عفونت راجعه}.
\end{itemize}
\item \lr{Scabies}: ناشی از مایت؛ بسیار خارش‌شونده.
\begin{latin}
Burrows in Interdigital spaces and flexor surfaces of the wrists, \textbf{elbow, axilla or genitals}
\end{latin}
\item هیستیوسیتوز سلول لانگرهانس: ضایعات پتشیال، هموراژیک.
\end{itemize}

\subsection*{هدف و اجزای درمان}
\begin{itemize}
\item هدف درمان در \lr{AD}: کاهش موارد \lr{flare} بیماری و طولانی نمودن مدت ریمیشن.
\item درمان: مرطوب کردن پوست؛ درمان دارویی برای خارش و التهاب؛ شناسایی و پرهیز از \lr{Ag}.
\end{itemize}

\clearpage
\PADheader{۳}{۸۰ ـ \lr{AD}}
\subsection*{مراقبت پوست}
\begin{itemize}
\item روزانه \lr{20 min} در وان حمام، سپس بلافاصله مرطوب‌کننده.
\item مرطوب‌کننده به صورت \textbf{پماد} (\lr{oint}) یا \textbf{کرم} (\lr{Cream}) باشد. \textbf{لوسیون مناسب نیست}.
\item بهتر است از پاک‌کننده‌های غیرصابونی استفاده شود.
\end{itemize}

\subsection*{درمان فاز حاد}
\begin{itemize}
\item داروهای کورتیکواستروئید و ایمونومدولاتور موضعی.
\item \textbf{پماد کورتیکواستروئید} بهتر است.
\item مهارکننده‌های کلسینورین (\textbf{تاکرولیموس}، پیمکرولیموس) (ایمونومدولاتور): خط دوم.
\item اکثر بیماران با مرطوب‌کننده و کورتیکواستروئید ضعیف قابل درمان است.
\item مهارکننده‌های کلسینورین: درمان کوتاه‌مدت و میان‌مدت در افراد \textbf{بالای ۲ سال} که به درمان‌های خط اول پاسخ نمی‌دهند.
\item در تمام مناطق بدن با عارضهٔ کم قابل استفاده است.
\item خط سوم درمان: \textbf{\lr{Crisaborole}} (مهارکنندهٔ فسفودی‌استراز ۴ (\lr{PDE-4})).
\item بیماری خفیف تا متوسط در بیماران \textbf{بالای ۲ سال} که به درمان‌های دیگر جواب نمی‌دهند.
\end{itemize}

\subsection*{درمان موارد شدید}
\begin{itemize}
\item درمان سیستمیک: کورتیکواستروئید، سیکلوسپورین، \lr{dupilumab}، مایکوفنولات موفتیل، متوترکسات.
\item اشعهٔ ماوراءبنفش.
\end{itemize}

\begin{itemize}
\item تا ۲۰\% موارد \lr{AD} به طور کامل بهبود پیدا می‌کنند.
\item لباس‌ها ۱۰۰\% کتان بهتر هستند.
\end{itemize}

\clearpage
\PADheader{۴}{۸۰ ـ \lr{AD}}
\subsection*{عوارض \lr{AD}}
\begin{itemize}
\item عفونت \lr{HSV} پوستی به صورت \textbf{سوپراینفکشن روی استاف طلایی}؛ اگزما هرپتیکوم؛ می‌تواند \textbf{تهدیدکننده} شود.
\item \lr{HPV}، وارت.
\item مولوسکوم کنتاژیوزوم.
\end{itemize}

\subsection*{الگوی تظاهرات بالینی \lr{AD}}
\begin{enumerate}
\item ضروری:
\begin{itemize}
\item \textbf{خارش}.
\item \textbf{اگزما صورت و سطوح اکستانسور} (در شیرخواران و کودکان).
\item \textbf{اگزما فلکسور} (در بزرگسالان).
\item درماتیت مزمن یا عودکننده.
\end{itemize}
\item شایع:
\begin{itemize}
\item سابقهٔ فردی یا خانوادگی، خشکی پوست، عفونت پوست.
\item \textbf{افزایش سطح \lr{IgE} سرم}.
\item درماتیت غیراختصاصی کف دست و پا.
\item سن پایین شروع، تست پوستی مثبت.
\end{itemize}
\end{enumerate}

\clearpage
\section*{پیوست: تصاویر منبع یادداشت‌های پیشین}
% Resolve paths when compiled from this directory or the repository root.
\IfFileExists{../../../dermatology/atopic dermatitis/images/pic01.jpg}
  {\def\PADoldimages{../../../dermatology/atopic dermatitis/images/}}
  {\def\PADoldimages{dermatology/atopic dermatitis/images/}}
\newcommand{\PADoldimage}[1]{\begin{center}
\textbf{تصویر پیشین \lr{#1}}\par\smallskip
\includegraphics[width=\textwidth,height=0.87\textheight,keepaspectratio]{\PADoldimages pic#1.jpg}
\end{center}\clearpage}
\PADoldimage{01}
\PADoldimage{02}
\PADoldimage{03}
\PADoldimage{04}
\PADoldimage{05}

\section*{پیوست: چهار تصویر جدید منبع}
\IfFileExists{images/pic01.jpg}
  {\def\PADnewimages{images/}}
  {\def\PADnewimages{pediatrics/ped-immunology/atopic dermatitis/images/}}
\newcommand{\PADnewimage}[1]{\begin{center}
\textbf{تصویر جدید \lr{#1}}\par\smallskip
\includegraphics[width=\textwidth,height=0.87\textheight,keepaspectratio]{\PADnewimages pic#1.jpg}
\end{center}\clearpage}
\PADnewimage{01}
\PADnewimage{02}
\PADnewimage{03}
\PADnewimage{04}
\endgroup
\PADfinish
